% ===== §14 Reception: The Ethics of Letting In =====
\section{The Ethics of Reception}
\label{p22:sec:reception}

\epigraph{\zh{至人之用心若镜，不将不迎，应而不藏。}}{\textit{Zhuangzi}}

\noindent
Of the four movements, reception is the one most easily passed over with a phrase. One says: do not defend, let it in. Yet reception is the step without precedent in the older forms of aesthetic education, and it is the step in which this paper's account is most its own. The reason the older forms lack it is instructive. Where beauty is a property of an object, there is nothing to receive; the masterwork stands there to be appraised, and the appraiser need not be altered by it. Reception becomes a distinct and demanding act only once beauty is relational, fragile, and passing, the registration of a coupling that gains, for such a beauty must be let into the one who perceives it or it does not survive the perceiving. To receive is to allow what has been seen to alter the valuation it enters, rather than to hold it at the level of cognition or to turn it back at the gate of defence. Stated so, reception looks passive. It is not. It is among the most strenuous things a relational subject does, and the traditions that have studied it disagree, as sharply as those that studied perception, about what kind of act it is.

\subsection{The semantic pivot: to receive beauty is to receive the other}
\label{p22:sub:reception-pivot}

Reception has two senses that must be separated before their unity can be claimed. In the first sense, what is received is the perceived beauty or appropriateness itself: the fragile, passing thing is let in, allowed to enter and to alter the one who perceives, rather than intercepted by defence or held off as mere information. This is reception as a matter of aesthetic psychology. In the second sense, what is received is the other: since the beauty of relational aesthetics is a coupling event, to receive the beauty is at root to receive the one with whom the beauty is jointly produced, together with that one's irreducibility, unpredictability, and refusal of the perceiver's control. This is reception as a matter of ethics.

The framework of this series does not let these two senses fall apart. Within relational aesthetics they are one act under two descriptions, because the beauty is produced precisely in the coupling with the other, so that to let the beauty in is to let the other in, and to defend against the other is to lose the beauty at its source. This is why an account of reception cannot remain in aesthetic psychology but slides of necessity into ethics, and the slide is not a digression. It is one more instance of the chain the whole paper has followed, the descent in which the aesthetic, pressed to its ground, discloses the ethical beneath it. Reception is the step where that descent is felt most directly, because here the thing one must let in has a face.

\begin{casebox}[The same remark, and what is done with it once it is seen]
Return to the morning and suppose the remark has been seen, registered as bearing something, a slight weight the routine did not predict. Seeing it is not yet receiving it. The hearer may, having seen, hold the remark off: note it as information, file it as the sort of thing she says, decide how to handle it, and in all of this keep it from entering and altering the morning, met at the level of management and not let in. Or the hearer may let it in: allow the remark to change, however slightly, the valuation of the morning it entered, let what was not foreseen revise what was planned, receive the small weight as a thing that may rewrite the next hour rather than as an obstacle to it. The difference is not in whether the remark was perceived, since both hearers perceived it, but in whether it was admitted, and the whole question of reception is the question of what divides the hearing that manages from the hearing that lets in. When the remark carries not delight but difficulty, an unforeseen need or a small reproach, the difference becomes the difference between defence and welcome, and the face in the remark is most plainly what is being kept out or let in.
\end{casebox}

\subsection{Levinas: the welcome of the other}
\label{p22:sub:reception-levinas}

The account that takes reception most seriously is the one that made the welcome of the other into first philosophy. On this account the other is encountered not first as an object of knowledge but as a face that breaks the perceiver's self-enclosed totality and addresses a claim that cannot be anticipated or absorbed into the categories already held. To welcome the other is to receive what exceeds every concept I would use to receive it, and the welcome is therefore not an act I perform from a prior freedom but a responsibility that precedes and constitutes my freedom~\parencite{levinas_totality}. The deepest temptation, on this account, is to reduce the other to the same, to take the face as one more item to be known and placed, and in that reduction the otherness of the other, which was the whole of what addressed me, is lost.

This gives the reception step its most exacting articulation. The injunction not to put the unexpected back in its place, which the four-step model stated almost casually, is here given its philosophical depth: to put back in place is precisely the reduction of the other to the same, the violence by which what exceeds my categories is forced to fit them. Reception, on this reading, is the refusal of that violence, the holding-open in which the other is let be as other, not assimilated and not mastered. The commitment is that the ethical precedes the ontological and that the other is absolutely prior, and the route gives the paper something it needs and could not get from psychology alone: an account of why reception is hard that locates the difficulty not in a deficit of the perceiver but in the structure of the encounter itself, in the irreducible excess of the one received over every means of receiving.

\subsection{Psychoanalysis: defence, containment, holding}
\label{p22:sub:reception-psychoanalytic}

If the welcome of the other says why reception is exacting, a second tradition says why it fails, and what would have to be developed for it to succeed. Here the obstacle to reception is named with clinical precision: the mechanisms of defence, denial and projection and isolation, are exactly the apparatus by which what is unbearable is kept from entering and altering the self. Against this stands a developmental account of the capacity that defence forecloses. In the figure of containment, one who is able takes in the raw, unmetabolised distress that another cannot yet bear, holds it, transforms it into a form that can be borne, and returns it~\parencite{bion_learning}; in the figure of the holding environment, a sustaining presence makes it safe enough for what is vulnerable to come forward~\parencite{winnicott_maturational} without being either annihilated or expelled. To receive, in this language, is to be able to contain what is not yet digested, what may disturb the one who receives it, without immediately defending against it or discharging it back out.

The commitment of this route is that reception is a psychological capacity, developed rather than given, and capable of being damaged. This matters to the paper in a way that is not merely illustrative. It supplies the mechanism behind the claim, made when the four steps were placed under the dialectic, that reception is the principal site after a relation has been wounded: injury raises defence, and raised defence forecloses the very containment that reception requires, so that what is perceived and even offered cannot get in. The route also introduces, against the most exalted reading of the welcome of the other, a needed note of mercy. A subject whose capacity to contain has been damaged by injury is not simply refusing the other; the gate is shut by a mechanism that the subject did not choose and cannot open by being commanded to. The tension between this and the ethical absoluteness of the welcome is real, and it is held for the section on conflicts rather than dissolved here.

\subsection{Gelassenheit: the posture of letting be}
\label{p22:sub:reception-gelassenheit}

A third figure describes not the ethics of reception nor its psychology but its posture, the bearing of the one who receives well. It is named as a releasement, a letting-be that is neither activity nor passivity in the ordinary sense, a comportment in which one steps out of the stance of representing-and-calculating, the stance that frames whatever it meets as a standing reserve to be ordered, and lets beings be the beings they are~\parencite{heidegger_gelassenheit}. Reception, under this figure, is the withdrawal from the posture of mastery, the loosening of the grip by which the will would seize and arrange what comes, so that the beauty and the other may arrive of themselves rather than be procured. The figure does not add a fourth mechanism so much as describe the disposition that the other routes presuppose: the welcome of the other and the capacity to contain both require that the one who receives has stepped out of the will to master, and releasement is the name for that step. It connects, as well, to the phenomenological route of the perception section, but its accent is different. There the letting-be served seeing, the coming-forward of the thing; here it serves keeping, the allowing of what has come forward to stay and to alter, without being seized.

A further phenomenological account refines the posture by describing what reception receives, and it is worth admitting because it names a structure the other accounts presuppose. On this account some phenomena give themselves, arrive as gifts that exceed what the receiver could have constituted from her own resources, saturating her capacity to receive rather than fitting neatly within it~\parencite{marion_givenness}; the face of the other, the event of beauty, are given in this way, arriving with more than the receiver brought, overflowing the concepts that would contain them. Reception, under this description, is the receiving of what gives itself in excess of the receiver, and its difficulty is the difficulty of receiving a gift too large for the hands held out, of letting in what cannot be reduced to what one already had. This sharpens the welcome of the other and the mirror alike: the other is received well when she is let arrive as the gift she is, in excess of the receiver's prior concepts, and badly when she is cut down to fit them, received only insofar as she matches what was expected. The reception step, on this account, is the holding-open of a capacity to receive more than one brought, a readiness for the given that does not pre-limit it to the constitutable, and the failure of reception is the reduction of the given to the receiver's prior measure, the refusal to be given more than one had. This is the same structure the welcome named as the refusal to reduce the other to the same, described now from the side of what is received rather than the side of the one who receives, and the two descriptions together fix the reception step as the admission of an excess, the letting-in of what arrives larger than the room prepared for it.

\subsection{The mirror that does not store: an Eastern principal voice}
\label{p22:sub:reception-mirror}

The principal Eastern voice on reception belongs to the same roots as the perception section's contemplative routes but speaks to a different moment, and it should not be confused with them: there the question was how to see, here it is how the heart may meet things without being snagged by them. The figure is the mirror. The heart of the realised person is used like a mirror, the Zhuangzi says, neither seeing things off nor going out to greet them, responding and yet not storing~\parencite{zhuangzi}, \zh{应而不藏}. The mirror receives all that comes before it, perfectly and without resistance, and retains nothing when it passes; the teaching of non-abiding in the Diamond Sutra, that one should give rise to the heart that dwells nowhere, names the same bearing from the other tradition. To receive, on this figure, is to meet what comes without defending against it and without grasping after it, neither pushing away nor holding on.

The value of this voice is that it names, in a single image, both of the ways reception fails, and they are not the same way. One failure is defence, the pushing-away, the refusal to let the beauty or the other in; against this the mirror does not see things off. The other failure is possession, the grasping-after, the seizing of the beauty as a thing to be had and hoarded; against this the mirror does not go out to greet, and responds yet does not store. The Western accounts tend to foreground the first failure, defence, and the apparatus against it. The mirror holds defence and possession in one figure and refuses both at once, and in doing so it supplies the term the section most needs, because the second failure, possession, is the one that connects reception forward to reproduction. A reception that grasps and stores is a reception that has begun to hoard, and the metaphysics of hoarding is exactly what the reproduction section must keep the good cycle clear of. The mirror, responding yet not storing, is the figure in which reception is already turned toward a reproduction that accumulates without possessing.

The figure repays being read clause by clause, because each of its three phrases names a distinct discipline of reception. Not seeing off: the mirror does not push away what comes before it, does not refuse the unwelcome image or distort the unflattering one, and the receiver who does not see off is the one who does not defend against the remark that carries difficulty, does not turn it back at the gate because it was unwelcome. Not going out to greet: the mirror does not reach toward what has not yet come, does not anticipate or solicit or project the image it would prefer, and the receiver who does not go out to greet is the one who does not seize the remark and make of it more than it was, does not grasp the beauty and clutch it as a possession. Responding yet not storing: the mirror answers perfectly to what is before it and keeps nothing when it passes, and the receiver who responds yet does not store is the one who lets the remark fully in, lets it alter the morning completely while it is present, and does not hoard it afterward as a held thing to be kept against loss. The three phrases are three refusals, of defence, of grasping, and of hoarding, and together they describe a reception that is total in the moment and possessive in no moment, which is exactly the reception relational aesthetics requires, a full admission of the other that seizes nothing.

Returned to the morning, the mirror sorts the failures the case left open. The hearer who manages the remark, holds it off as information to be handled, is seeing it off, defending against an image she will not let alter her; she fails by the first refusal. The hearer who seizes the remark, makes of a small unforeseen weight a large significance she clutches, who takes the difficult remark as a wound to be nursed or the warm one as a possession to be hoarded, is going out to greet and storing; she fails by the second and third. The mirror's receiver does neither: she lets the remark in fully, lets it revise the morning while it is present, and keeps nothing of it as a held grievance or a hoarded treasure, so that when the next moment comes she meets it as freshly as she met this one, undistorted by what she is carrying from before. This is the bearing that the reception step cultivates, and the case shows it to be a narrow path between two failures that the Western accounts, fixed on defence alone, do not jointly see.

\subsection{Two real conflicts}
\label{p22:sub:reception-conflicts}

The accounts of reception conflict at two points, and as with perception the framework resolves them by taking sides rather than by averaging.

The first conflict sets the ethical absoluteness of the welcome against the developmental contingency of the capacity to contain. On the one reading, to receive the other is an unconditioned demand, prior to my freedom and admitting no excuse; to plead that I lack the capacity is very nearly to evade the responsibility that constitutes me. On the other reading, the capacity to contain is a fragile achievement that injury can damage, and to issue an unconditioned demand for reception to a subject whose containment has been wounded is not exalted but cruel, since it commands an opening that the mechanism of defence has shut and that no command can reopen. The conflict is real and it carries a practical stake the paper will face directly when it asks, in the political economy of the field, what may be required of those whom relentless bad repetition has worn down. The resolution is not to choose one description and discard the other but to assign them to different offices. The welcome of the other states the norm of reception, what reception is when it succeeds, the direction in which the good lies; it is the account of the end. The developmental capacity states the condition of reception, what must be intact or repaired for the norm to be met; it is the account of the ground that the end requires. The series has held this couple before, the end that orients and the condition that constrains, and has insisted that an ethic which protects the conditions of reception is more just than one which demands the end while the conditions lie in ruins. Mercy toward the wounded gate is not a relaxation of the norm. It is the recognition that the norm is reached by way of the condition, and that to attend to the condition is to serve the norm rather than to excuse its abandonment.

The assignment to different offices can seem too neat, and a sharper form of the conflict resists it. Suppose the wounded subject's defence is not a damaged capacity awaiting repair but a settled refusal dressed as incapacity, the gate held shut by a will that finds in the language of woundedness a permission never to open. The developmental account, pressed this way, threatens to supply an alibi: any failure to receive can be redescribed as an injury that has foreclosed containment, and the unconditioned demand, which alone would expose the refusal, is silenced in advance by the plea of incapacity. If every closed gate is a wounded one, the ethical absoluteness of the welcome has no purchase, and reception becomes optional for anyone willing to call their refusal an injury. This is a real danger and the framework does not escape it by fiat. What it offers is a mark of the difference, drawn from the framework's own terms: a genuinely wounded capacity shows the signature of a capacity under repair, a gate that lowers as the injury heals, a reception that returns as the field is tended, whereas a refusal dressed as incapacity shows no such trajectory, the gate staying shut regardless of how the conditions improve, the incapacity invoked anew at each occasion without ever being worked through. The difference is legible over time, in whether the condition, attended to, yields a reception that grows, or whether the plea of incapacity is a fixed point that no improvement of conditions moves. Mercy toward the wounded gate is owed to the gate that is healing; it is not owed indefinitely to a gate that uses the name of woundedness to stay shut, and the framework's insistence on reading the field over time, rather than accepting a single moment's plea, is what keeps the developmental account from becoming the alibi it could otherwise supply.

The second conflict is the dangerous one, and it is the structural twin of the hardest step in the perception section. The mirror responds yet does not store; the Daoist heart meets things without abiding in them. Yet the fourth movement of the model requires that beauty be \emph{reproduced}, that the relational memory thicken, that the threshold of joint perception be lowered by what has gone before. To thicken, to lower a threshold by accumulation, looks exactly like storing, like the abiding that the mirror forbids. If reception must not store, and reproduction must accumulate, the model appears to demand at one step what it prohibits at another.

\subsection{The hardest step: responding-yet-not-storing grounds healthy reproduction}
\label{p22:sub:reception-hardest}

The resolution turns on what kind of thing is accumulated, and it is the same distinction the series has drawn between value that is held as a possessed stock and value that is the changed curvature of a field. What the mirror forbids is possessive retention, the seizing of the beauty or the other as a thing to be had, kept against loss, hoarded as a private holding. This is the imaginary register of value, the idolised possession whose nature is to be clung to and whose clinging exhausts what it clings to. The accumulation that healthy reproduction requires is not of this kind. It is not a stock laid up somewhere but the raising of the next cycle's starting point, the changed disposition of the field, the lowered threshold at which the partners can jointly perceive what before went unseen. Nothing is held; the curvature has changed. Holonomy accumulates not as a thing retained but as a property of the path, and a property of the path is precisely what cannot be stored, because it is not the sort of thing that occupies a place from which it could be kept or lost.

\claim{Responding-yet-not-storing grounds healthy reproduction rather than forbidding it. What the mirror prohibits is possessive retention, the hoarding of beauty or the other as a held stock, which belongs to the imaginary register and exhausts what it grasps. What reproduction accumulates is not a stock but a changed curvature of the field, a raised starting point for the next cycle, which is a property of the path and not a thing in a place. Because such accumulation is not storing, it is the only kind of accumulation a non-grasping reception can sustain; the mirror's refusal to store is therefore the condition of an accumulation that does not exhaust, not an obstacle to it.}

So the second conflict, like the third conflict of the perception section, inverts. The mirror's refusal to abide does not forbid the thickening that reproduction needs; it founds the only thickening that can go on without exhausting its source, the accumulation that gains because it does not possess. This is the ancient counsel once more, to generate without possessing and to bring to maturity without lording over, read now into the conditions of reception and carried forward into time. A reception that grasped would feed a reproduction that hoarded, and a reproduction that hoards is the bad cycle, the spiral collapsed into a clutching wheel. A reception that responds and does not store feeds a reproduction that accumulates as curvature, and that is the good cycle continuing itself. The mirror, it turns out, was already a teaching about time.

\subsection{A layered synthesis}
\label{p22:sub:reception-synthesis}

The accounts of reception take their places without identification, each at the level where it is true. The welcome of the other is taken at the level of ethics: it states the norm of reception, the direction of the good, the refusal to reduce the other to the same, and it supplies the reason reception is exacting that no psychology could supply, the irreducible excess of the received over the means of receiving. The figures of containment and holding are taken at the level of implementation and development: they describe reception as a capacity that can be built and damaged, they ground the claim that injury makes reception the principal site, and they may be read alongside an observable account of how unexpected events are integrated rather than defended against. The mirror that responds yet does not store is taken at the level of contemplative cultivation and of metaphysical guard: it names both failures reception must avoid, defence and possession, and it holds the non-possessive line that keeps reception from feeding a hoarding reproduction, and through its internal distinction between retaining a stock and changing a curvature it secures the accumulation the framework needs. Releasement runs through all three as the posture they presuppose, the withdrawal from mastery without which neither welcome nor containment nor non-abiding is possible.

The arrangement uses no equals sign, because the welcome of the other is not the capacity to contain, and neither is the mirror, and to write them as equal would be to claim an identity that none of them would accept and that the framework does not need. What is claimed instead is an arrangement in which each is preserved as itself and asked only for what it can give, which is once more the seeking of common ground while the differences are kept. Reception, gathered this way, is shown to be the step at which the aesthetic discloses the ethical most plainly, because the thing that must be let in is not a sensation but a face, and the discipline of letting it in without defence and without seizure is, at one and the same time, the discipline of a beauty that is produced only between.

\subsection{The reach and the defeat of each account}
\label{p22:sub:reception-ledger}

The arrangement of reception's accounts, like that of perception's, should rest on a frank accounting of what each grasps and what defeats it, rather than on a tidy distribution of levels, and the accounting is owed especially here because the accounts of reception differ not only in what they explain but in what they take reception fundamentally to be, an ethical relation, a psychological capacity, a contemplative bearing.

The ethics of the welcome explains what no psychology can, why reception is exacting in principle and not merely difficult in fact, locating the difficulty in the irreducible excess of the other over every means of receiving her, and it supplies the norm without which reception would have no direction, the refusal to reduce the other to the same. What defeats it, taken alone, is its silence about capacity and its consequent harshness. By making reception an unconditioned demand prior to the subject's freedom, it has no place for the subject whose capacity to receive has been damaged, and pressed to the letter it convicts the wounded of evasion when they plead an incapacity the welcome cannot acknowledge. Its strength is the norm and its defeat is the absence of any account of the conditions under which the norm can be met, which is why it cannot be the whole and must be joined to an account of capacity it cannot itself provide.

The psychoanalytic account explains exactly what the welcome cannot, the capacity to receive as something developed and damageable, and it grounds the paper's most humane claims, that injury raises defence and that reception can be cultivated and repaired. It is the account that can be operationalized, that connects reception to observable integration and gating, and that gives the wounded gate a mechanism and a path to healing. What defeats it, taken alone, is its silence about the norm. It describes the capacity and its vicissitudes without telling us what reception is for, toward what it is exercised, and a complete account of reception cannot rest on a description of capacity that is silent about the good the capacity serves. Its strength is the mechanism of capacity and its defeat is the absence of the norm that the welcome supplies, the two being complementary in exactly the way that a level-account requires, the one giving the end and the other the condition.

The contemplative figure of the mirror explains what neither of the others foregrounds, that reception fails in two opposite ways and not one, by defence and by possession, and it holds the non-possessive line that keeps reception from feeding a hoarding reproduction. What defeats it, taken as a complete account rather than as a discipline and a guard, is the same limit that defeated the contemplative route in perception, its tendency at the metaphysical limit to a non-attachment so complete that it would empty reception of the investment the paper requires, the caring that the non-abiding heart must still, as the perception section argued, sustain. The mirror is honoured as the discipline that names both failures and the guard against possession, and it is not adopted as a final doctrine of non-attachment, because a reception wholly without abiding would be a reception without the caring investment that distinguishes it from indifference.

The materialist resolution here is the same in form as in perception and different in content. It refuses to crown any of the three, the ethics of the welcome, the psychology of capacity, the discipline of the mirror, as the whole truth of reception, holding instead that reception is a real relation with a real other in a real field, possessing at once a norm, a condition, and a discipline, none of which is reducible to the others. It is materialist in taking the other received to be a real other, irreducible and unconstructed, and the capacity to receive her to be a real and damageable feature of a material subject, not a transcendental structure nor a moment of spirit; it is dialectical in holding that the norm, the condition, and the discipline are aspects of one real relation whose truth is their articulation and not any one of them. Against the idealism that would make the other a moment in the subject's self-relation, the paper holds the other's material irreducibility; against the reductive psychology that would make reception merely a regulation of arousal, it holds the norm that the regulation serves; against the contemplative absolutism that would make non-attachment the whole, it holds the caring investment that reception requires. Reception is a real relation to a real other, normed and conditioned and disciplined, and reducible to none of its aspects, which is the materialist dialectic applied to the ethics of letting in.
